\chapter{\ifenglish Introduction\else บทนำ\fi}\label{ch:intro}

\section{\ifenglish Project rationale\else ที่มาและความสำคัญของโครงงาน\fi}

การเรียนรู้ในระดับอุดมศึกษาเกิดขึ้นภายในระบบนิเวศที่ประกอบด้วยผู้มีส่วนร่วมหลายฝ่าย ได้แก่ ผู้เรียน ซึ่งเป็นผู้สร้างความรู้และพัฒนาทักษะของตนเอง อาจารย์ ซึ่งเป็นผู้ออกแบบประสบการณ์การเรียนรู้และประเมินผล และเพื่อนร่วมชั้น ซึ่งเป็นทั้งผู้ร่วมงานและผู้สะท้อนพฤติกรรมการทำงานของกันและกัน สิ่งที่เชื่อมโยงทุกฝ่ายเข้าด้วยกันคือ ``กระบวนการการให้ข้อเสนอแนะ'' ซึ่งเป็นกลไกที่ทำให้ผู้เรียนรู้ว่าตนเองอยู่ตรงไหน ควรปรับปรุงอะไร และจะไปต่อทางใด

ข้อเสนอแนะที่มีคุณภาพไม่ได้จบที่การให้คะแนน แต่ต้องทำให้ผู้เรียนนำไปพัฒนาตนเองได้จริง กล่าวคือ ต้องเฉพาะเจาะจง ทันเวลา และเชื่อมโยงกับเป้าหมายการเรียนรู้ อย่างไรก็ตาม ในทางปฏิบัติ ข้อเสนอแนะลักษณะนี้ต้องอาศัยเวลาและแรงงานของอาจารย์เป็นอย่างมาก โดยเฉพาะในรายวิชาที่มีนักศึกษาจำนวนมาก (large-section courses) และมีการทำงานเป็นกลุ่ม อาจารย์มักไม่สามารถติดตามการมีส่วนร่วมของนักศึกษาแต่ละคนได้อย่างทั่วถึง ทำให้ข้อเสนอแนะที่นักศึกษาได้รับมีลักษณะกว้างและเหมือนกันทั้งกลุ่ม ไม่สะท้อนบทบาทและการพัฒนารายบุคคล นอกจากนี้การประเมินตนเองและประเมินเพื่อนมักทำเป็นครั้งคราว หรือทำเพียงเพื่อนำคะแนนไปคำนวณ จึงไม่เกิดการสะท้อนคิด (reflection) อย่างต่อเนื่อง ในด้านภาระงานของอาจารย์ ต้องอ่านผลประเมินและข้อความจำนวนมากด้วยตนเอง ทั้งเพื่อสรุปภาพรวม คัดแยกประเด็นที่ต้องติดตาม และเขียนข้อเสนอแนะตอบกลับ ซึ่งใช้เวลามากและทำได้ไม่ครบทุกคน ยิ่งไปกว่านั้น ในการประเมินเพื่อนอาจมีข้อความที่รุนแรง ไม่สร้างสรรค์ หรือเป็นการกล่าวโจมตีตัวบุคคล ซึ่งอาจส่งผลเสียต่อความสัมพันธ์ในกลุ่มและสุขภาวะของนักศึกษา แต่อาจารย์อาจตรวจพบไม่ทัน

ด้วยเหตุนี้เพื่อแก้ไขปัญหาดังกล่าว คณะผู้จัดทำจึงพัฒนา Collab Reflection แพลตฟอร์มที่ให้นักศึกษาประเมินตนเองและประเมินเพื่อนในกลุ่มอย่างเป็นระบบ โดยมุ่งเน้นให้ ``ข้อเสนอแนะเป็นเครื่องมือพัฒนาผู้เรียน'' อย่างเต็มรูปแบบ และใช้ปัญญาประดิษฐ์ (AI) เป็นผู้ช่วยลดภาระงานของอาจารย์ เพื่อให้อาจารย์เข้าถึงนักศึกษาจำนวนมากได้โดยไม่ลดคุณภาพของข้อเสนอแนะ ทั้งนี้ AI ทำหน้าที่เป็นผู้ช่วยเท่านั้น ไม่ได้แทนที่การตัดสินใจของอาจารย์ โดยอาจารย์ยังคงเป็นผู้ตรวจทานและเป็นผู้ให้ข้อเสนอแนะขั้นสุดท้ายแก่นักศึกษา

\section{\ifenglish Objectives\else วัตถุประสงค์ของโครงงาน\fi}

\begin{enumerate}
  \item เพื่อพัฒนาเว็บแอปพลิเคชัน Collab Reflection ที่ประกอบด้วยแบบประเมินตนเอง แบบประเมินเพื่อนในกลุ่ม และแดชบอร์ดสำหรับอาจารย์ผู้สอนที่แสดงและจำแนกประเภทผลการประเมิน
  \item เพื่อพัฒนาระบบ AI ที่ตรวจโทนข้อความในการประเมินเพื่อน และเสนอถ้อยคำเชิงสร้างสรรค์ให้ผู้เขียนเลือกใช้
  \item เพื่อพัฒนาระบบปัญญาประดิษฐ์ (AI) ที่สรุปผลการประเมิน เพื่อช่วยอาจารย์ร่างข้อเสนอแนะตอบกลับนักศึกษา โดยอาจารย์เป็นผู้ตรวจทานก่อนส่งทุกครั้ง
  \item เพื่อลดภาระของอาจารย์ในการให้ข้อเสนอแนะ และให้นักศึกษาได้รับข้อเสนอแนะที่เหมาะสมและทั่วถึง สำหรับนำไปพัฒนาการเรียนรู้และการทำงานร่วมกับผู้อื่น
\end{enumerate}

\section{\ifenglish Project scope\else ขอบเขตของโครงงาน\fi}

ผู้ใช้งาน : นักศึกษาและอาจารย์ผู้สอน

\subsection{ขอบเขตด้านฟังก์ชัน}

\begin{itemize}
  \item นักศึกษา
  \begin{enumerate}
    \item แบบประเมินตนเอง
    \item แบบประเมินเพื่อนในกลุ่ม
    \item ปัญญาประดิษฐ์ (AI) ตรวจโทนข้อความในการประเมินเพื่อน และเสนอถ้อยคำเชิงสร้างสรรค์ให้ผู้เขียนเลือกใช้
  \end{enumerate}
  \item อาจารย์ผู้สอน
  \begin{enumerate}
    \setcounter{enumii}{3}
    \item แดชบอร์ดสำหรับดูและจำแนกประเภทผลการประเมิน
    \item ปัญญาประดิษฐ์ (AI) สรุปผลการประเมินและช่วยร่างข้อเสนอแนะตอบกลับนักศึกษา โดยอาจารย์เป็นผู้ตรวจทานก่อนส่งทุกครั้ง
  \end{enumerate}
\end{itemize}

\subsection{ขอบเขตด้านการพัฒนา}

โครงงานแบ่งการพัฒนาเป็น 2 ระยะ ปัจจุบันเสร็จสิ้นการออกแบบระบบและการสาธิตต้นแบบ (demo) แล้ว และกำลังพัฒนาเว็บแอปพลิเคชันในระยะที่ 1 ซึ่งครอบคลุมแบบประเมินตนเอง แบบประเมินเพื่อนในกลุ่ม และแดชบอร์ดสำหรับอาจารย์ ส่วนระบบปัญญาประดิษฐ์ (AI) ทั้งสองส่วน (ตรวจโทนข้อความและสรุปผลร่างข้อเสนอแนะ) จะพัฒนาในระยะที่ 2

\subsection{ข้อจำกัดของขอบเขต}

\begin{itemize}
  \item ระยะที่ 1 ยังไม่รวมระบบปัญญาประดิษฐ์ (AI)
  \item ระบบยังไม่ได้ผสานกับระบบของมหาวิทยาลัย
  \item ในขั้นการออกแบบระบบและการสาธิตต้นแบบยังไม่ได้นำไปใช้งานในรายวิชาจริง
\end{itemize}

\section{\ifenglish Expected outcomes\else ประโยชน์ที่คาดว่าจะได้รับ\fi}

\begin{itemize}
  \item นักศึกษา
  \begin{itemize}
    \item ได้รับโอกาสสะท้อนคิดเกี่ยวกับการทำงานของตนเองและของกลุ่มอย่างสม่ำเสมอ
    \item ได้รับข้อเสนอแนะที่เฉพาะเจาะจงและเหมาะสมมากขึ้น ทำให้เห็นจุดแข็งและจุดที่ควรพัฒนาของตนเองได้ชัดเจน
    \item ได้ฝึกการให้ข้อเสนอแนะแก่เพื่อนด้วยถ้อยคำเชิงสร้างสรรค์ ซึ่งเป็นทักษะสำคัญในการทำงานร่วมกับผู้อื่น
  \end{itemize}
  \item อาจารย์ผู้สอน
  \begin{itemize}
    \item ประหยัดเวลาในการอ่านผลประเมิน คัดแยกประเด็น สรุปผล และเขียนข้อเสนอแนะ
    \item มองเห็นภาพรวมของผลการประเมินทั้งรายวิชาได้ง่ายผ่านแดชบอร์ด
    \item สามารถให้ความสำคัญกับนักศึกษาที่ต้องการความช่วยเหลือจริงได้ทันท่วงที
    \item ยังคงเป็นผู้ตรวจทานข้อเสนอแนะก่อนส่งทุกครั้ง จึงควบคุมคุณภาพของข้อเสนอแนะได้
  \end{itemize}
  \item รายวิชาและหน่วยงาน
  \begin{itemize}
    \item สามารถขยายการให้ข้อเสนอแนะเชิงพัฒนาไปยังนักศึกษาจำนวนมากได้อย่างทั่วถึง
    \item มีข้อมูลการประเมินที่เป็นระบบและจำแนกประเภทแล้ว สำหรับนำไปปรับปรุงการจัดกิจกรรมกลุ่มในรายวิชา
    \item เป็นแนวทางต้นแบบของการนำปัญญาประดิษฐ์ (AI) มาช่วยงานด้านการให้ข้อเสนอแนะในการเรียนการสอน โดยมีมนุษย์เป็นผู้ตัดสินใจขั้นสุดท้าย
  \end{itemize}
\end{itemize}

\section{\ifenglish Technology and tools\else เทคโนโลยีและเครื่องมือที่ใช้\fi}

\subsection{\ifenglish Hardware technology\else เทคโนโลยีด้านฮาร์ดแวร์\fi}

\begin{itemize}
  \item เครื่องคอมพิวเตอร์ส่วนบุคคลสำหรับพัฒนาระบบ
  \item เซิร์ฟเวอร์สำหรับติดตั้งและให้บริการระบบ
\end{itemize}

\subsection{\ifenglish Software technology\else เทคโนโลยีด้านซอฟต์แวร์\fi}

\begin{itemize}
  \item \textbf{React 19 (TypeScript) + Vite} -- พัฒนาส่วนหน้า (Frontend) ของเว็บแอปพลิเคชัน
  \item \textbf{Tailwind CSS} -- จัดรูปแบบและออกแบบหน้าตาของส่วนติดต่อผู้ใช้
  \item \textbf{TanStack Query และ Zustand} -- จัดการข้อมูลจากเซิร์ฟเวอร์และสถานะของแอปพลิเคชันฝั่งหน้าเว็บ
  \item \textbf{Node.js และ Express 5} -- พัฒนาส่วนหลังบ้าน (Backend) และ API
  \item \textbf{PostgreSQL และ Drizzle ORM} -- จัดเก็บข้อมูลการประเมิน และจัดการการเชื่อมต่อฐานข้อมูล
  \item \textbf{Qwen} -- โมเดลภาษา (LLM) สำหรับตรวจโทนข้อความ เสนอถ้อยคำเชิงสร้างสรรค์ และสรุปผลเพื่อช่วยอาจารย์ร่างข้อเสนอแนะ (ระยะที่ 2)
  \item \textbf{Docker และ nginx} -- จัดการสภาพแวดล้อมและติดตั้งระบบ (Deployment)
  \item \textbf{Vitest และ Cypress} -- ทดสอบระบบ (Unit Test และ End-to-End Test)
  \item \textbf{Git/GitHub และ Visual Studio Code} -- จัดการซอร์สโค้ดและพัฒนาโปรแกรม
\end{itemize}

\section{\ifenglish Project plan\else แผนการดำเนินงาน\fi}

แผนการดำเนินงานของโครงงานแสดงดังตารางที่~\ref{tab:ch1-plan} ณ วันที่จัดทำรายงานนี้ ขั้นตอนการค้นหางานวิจัยและกำหนดปัญหา การสัมภาษณ์อาจารย์ การออกแบบระบบและสร้างต้นแบบ (demo) รวมถึงการนำเสนอความก้าวหน้ารอบที่ 1 และ 2 ได้ดำเนินการเสร็จสิ้นแล้ว ส่วนการพัฒนาเว็บแอปพลิเคชันระยะที่ 1 อยู่ระหว่างดำเนินการ

% แผนภูมิแกนต์ใช้ environment plan ของ cpecmu.cls (\planitem{เดือนเริ่ม}{ปีเริ่ม ค.ศ.}{เดือนสิ้นสุด}{ปีสิ้นสุด}{ขั้นตอน})
% plan ไม่มีช่อง caption จึงใส่ด้วย \captionof (แพ็กเกจ capt-of ในไฟล์หลัก)
\captionof{table}{แผนการดำเนินงาน}\label{tab:ch1-plan}
\nopagebreak
\begin{plan}{7}{2026}{3}{2027}
    \planitem{7}{2026}{9}{2026}{ค้นหางานวิจัยและกำหนดปัญหา}
    \planitem{8}{2026}{8}{2026}{นำเสนอความก้าวหน้ารอบที่ 1}
    \planitem{9}{2026}{9}{2026}{สัมภาษณ์อาจารย์เพื่อรวบรวมความต้องการ}
    \planitem{8}{2026}{10}{2026}{ออกแบบระบบและสร้างต้นแบบ (demo)}
    \planitem{10}{2026}{10}{2026}{นำเสนอความก้าวหน้ารอบที่ 2}
    \planitem{11}{2026}{11}{2026}{รวบรวมความต้องการจากฝั่งนักศึกษา}
    \planitem{11}{2026}{12}{2026}{พัฒนาเว็บแอปพลิเคชันระยะที่ 1}
    \planitem{1}{2027}{2}{2027}{พัฒนาระบบปัญญาประดิษฐ์ (AI) (ระยะที่ 2)}
    \planitem{1}{2027}{2}{2027}{ทดสอบระบบ}
    \planitem{1}{2027}{3}{2027}{ทดลองใช้งานจริงกับรายวิชา}
    \planitem{2}{2027}{3}{2027}{ประเมินผลและสรุปผลการดำเนินงาน}
    \planitem{2}{2027}{3}{2027}{จัดทำรูปเล่มรายงาน}
\end{plan}

\section{\ifenglish Roles and responsibilities\else บทบาทและความรับผิดชอบ\fi}

คณะผู้จัดทำแบ่งบทบาทและหน้าที่ตามความถนัด ดังนี้

\begin{itemize}
  \item นายปฏิพันธ์ เลขนอก (Project Manager \& Backend) -- วางแผนและติดตามความคืบหน้าของโครงงาน ประสานงานระหว่างสมาชิกและอาจารย์ที่ปรึกษา พัฒนาส่วน Backend API และออกแบบฐานข้อมูล ใช้ความรู้ด้านการบริหารโครงการ Software Engineering และการออกแบบระบบฐานข้อมูล
  \item นางสาวดารารัตน์ บัวยิ้ม (Full-Stack \& Validation) -- พัฒนาส่วนหน้าเว็บของแบบประเมินและแดชบอร์ดอาจารย์ รวมถึงพัฒนา Backend API และเชื่อมต่อกับระบบปัญญาประดิษฐ์ (Qwen) ใช้ความรู้ด้าน Web Development, API Design และการประยุกต์ใช้โมเดลภาษา
  \item นายรัชชานนท์ ไชยวงค์ (Backend \& Tester) -- พัฒนาส่วน Backend และฐานข้อมูล วางแผนและดำเนินการทดสอบระบบทั้ง Unit Test และ End-to-End Test เพื่อตรวจสอบความถูกต้องของการทำงาน ใช้ความรู้ด้าน Software Testing, ฐานข้อมูล และการเขียน API
  \item นางสาวบัวชมพู ฤกษ์สุทธิรัตน์ (Frontend \& UX/UI) -- ออกแบบประสบการณ์และหน้าตาของผู้ใช้ (UX/UI) พัฒนาส่วนหน้าเว็บให้ใช้งานง่ายสำหรับนักศึกษาและอาจารย์ ใช้ความรู้ด้าน UX/UI Design, Frontend Development และการออกแบบส่วนติดต่อผู้ใช้
\end{itemize}
